\documentclass[10pt,a4paper]{amsart}
\usepackage[margin=19mm]{geometry}
\usepackage{amsmath,amssymb,mathtools}
\usepackage[scr=rsfs,cal=euler]{mathalfa}
\usepackage{xfrac}
\usepackage[unicode,hidelinks,bookmarks=false]{hyperref}
\newcommand{\ie}{i.e.\ }
\newcommand{\cf}{cf.\ }
\newcommand{\resp}{respectively\ }
\newcommand{\Subsection}{\subsection}
\providecommand{\nobreakdash}{\nobreak\hbox{-}}
\setlength{\emergencystretch}{2em}
\multlinegap0pt
\usepackage{bm}
\usepackage{bbm}
\usepackage{dsfont}
\usepackage{xspace}
\usepackage{cancel}
\usepackage{mathtools}
\usepackage{amsthm}
\makeatletter
\@ifundefined{theorem}{\newtheorem{theorem}{Theorem}}{}
\@ifundefined{claim}{\newtheorem{claim}[theorem]{Claim}}{}
\makeatother
\DeclareMathOperator{\dist}{d}
\DeclareMathOperator{\length}{length}
\newcommand{\Z}{\mathbf{Z}}
\newcommand{\into}{\hookrightarrow}
\title[Vertex-transitive graphs with uniformly bisecting quasi-geod]{Vertex-transitive graphs with uniformly bisecting quasi-geodesics}
\author{Joseph Paul MacManus}
\date{Source: arXiv:2511.10759v2 (CC BY 4.0)}
\begin{document}
\maketitle

\noindent\emph{Source and licence.} Vertex-transitive graphs with uniformly bisecting quasi-geodesics. Version arXiv:2511.10759v2 (CC BY 4.0), distributed under the Creative Commons Attribution 4.0 International licence, as recorded in the arXiv licence metadata for this exact version and shown on the abstract page https://arxiv.org/abs/2511.10759v2. Full rights notice: licenses/arxiv-2511.10759-CC-BY-4.0.txt.\par\smallskip

\noindent\textit{Editorial scope.} Counted unit: the Claim whose source label is \texttt{claim:quasi-circle} in the article (source0.tex lines 1597--1599, source label \texttt{claim:quasi-circle}), delivered together with the article's own complete original proof of it (source0.tex lines 1601--1653, one \texttt{proof} environment of 53 lines). No mathematics is written, repaired or re-derived: statement and proof are the article's own text line for line. Required context retained uncounted: eq:t-i property (lines 1501--1503); claim:bound-on-pi-dist (lines 1544--1546); claim:bound-on-dist-from-gamma-to-Tn (lines 1559--1562). Every definition, display and label of those lines is retained unchanged. Omitted from this package and not redistributed: 1--1500, 1504--1543, 1547--1558, 1563--1596, 1600--1600, 1654--1916. Nothing inside a retained excerpt is omitted or altered: every retained body is delivered exactly as the article's lines are written, so no omission receipt of any kind accompanies this package. Citations: the retained excerpts carry no citation command at all, so no bibliography entry of the article is transcribed and no reference is redistributed. Reference adaptation: none was needed and none was made; every reference inside the delivered unit resolves inside it, so no unresolved reference or citation is produced. Printed slips kept literal: no printed word, symbol, hypothesis or number of the counted statement or of its proof was altered. Layout and class: the article's own class is \texttt{amsart} with packages including inputenc, fontenc, amsfonts, amsmath, amsthm, amssymb, thmtools, caption, stmaryrd, mathrsfs, todonotes, graphicx, tikz, quiver; the wrapper is the lane's pinned \texttt{amsart} editorial wrapper at 10pt on A4, which restates the theorem-like environments the retained text needs and adds the article's own macro definitions that the retained text needs (\texttt{\textbackslash Z}, \texttt{\textbackslash dist}, \texttt{\textbackslash into}, \texttt{\textbackslash length}), with extra wrapper packages (bm, bbm, dsfont, xspace, cancel, mathtools). The delivered numbering is the unit's own. Assets: no retained span contains a figure environment, a graphics-inclusion command, a PDF-page inclusion, a table or a TikZ picture; no image file is copied into this package, the package declares no asset dependency and no image file is redistributed with it. Licence: version arXiv:2511.10759v2 of this article is distributed under the Creative Commons Attribution 4.0 International licence, as recorded in the arXiv licence metadata for this exact version and shown on the abstract page https://arxiv.org/abs/2511.10759v2. A full rights notice is delivered at licenses/arxiv-2511.10759-CC-BY-4.0.txt. Characters: every retained excerpt is pure ASCII, so the delivered engine, XeLaTeX with its default Latin Modern text font, reports no missing character.
    \begin{equation}\label{eq:t-i property}
        \frac{t_i}{\dist_X(\gamma_i(t_i), T_n)} \geq 2\lambda, 
    \end{equation}

    \begin{claim}\label{claim:bound-on-pi-dist}
        Given distinct $i , j \in \Z/3\Z$, we have that $\dist_X(p_i,p_j) \geq \frac 2{2\lambda+1} d_n - c$. 
    \end{claim}

    \begin{claim}\label{claim:bound-on-dist-from-gamma-to-Tn}
        For all $i \in \Z/3\Z$, we have that 
        $\dist_X(\gamma_i', T_n) > \frac{d_n}{2\lambda+1}$.
    \end{claim}

    \begin{claim}\label{claim:quasi-circle}
        We have that $Q$ is a $(\lambda', c')$-quasi-circle, where $\lambda' = 48\lambda^3$ and $c' = 2c$. 
    \end{claim}

    \begin{proof}
        We must show that the inclusion $Q \into X$ is a quasi-isometric embedding with uniform constants not depending on $n$. 
        
        Let $x,y \in Q$. If $x$, $y$ lie in a common `side' of $Q$ then we are trivially done. Suppose then that $x \in p_1$, $y \in p_2$. Then by Claim~\ref{claim:bound-on-pi-dist} and the earlier bound on $\length(Q)$, we have that 
        $$
        \dist_X(x,y) \geq \frac{1}{8\lambda^2(2\lambda+1)} \dist_Q(x,y) - 2c \geq \frac 1 {24\lambda^3} \dist_Q(x,y) - 2c.
        $$
        Similarly, if $x \in \gamma_1' \cup \gamma_2'$ and $y \in \alpha \subset T_n$, then by Claim~\ref{claim:bound-on-dist-from-gamma-to-Tn} we may compute
        $$
        \dist_X(x,y) \geq \frac 1 {48\lambda^3} \dist_Q(x,y) - c.
        $$
        Now, suppose without loss of generality $y \in \gamma_1' \cup \gamma_2'$ and $x \in p_1$, the case where $x \in p_2$ is identical. 
        If $y \in \gamma_2'$, then we can easily compute 
        \begin{align*}
            \dist_X(x,y) &\geq \frac 1 \lambda t_1 - c - \length(p_1) \\
            &\geq \frac{d_n}{2\lambda+1} - c \\
            & \geq \frac{1}{48\lambda^3} \dist_Q(x,y) - 2c.
        \end{align*}
        Suppose instead that $y \in \gamma_1'$. To ease calculation, let us reverse and re-parametrise $\gamma_1'$ so that $\gamma_1'(0) = p_1(0)$. Then, write $x = p_1(a)$, $y = \gamma_1'(b)$. Let $D = \dist_X(x,y)$. 
        Suppose for the sake of a contradiction that $D < \frac 1 {8\lambda} (a+b) - c$.
        By our choice of $t_1$ earlier as the minimal argument satisfying (\ref{eq:t-i property}), we have that 
        $$
            t_1 - b < 2\lambda \dist_X(y, T_n) \leq 2\lambda (\dist_X(x,y) + \dist_X(x,T_n)) = 2\lambda(D + \length(p_1) - a).
        $$
        As $\length(p_1) \leq \frac {t_1}{2\lambda}$, we see that
        $t_1 - b < 2\lambda D - 2\lambda a + t_1$ , and so $2\lambda a < b + 2\lambda D$. 
        Applying our hypothesised upper bound on $D$, routine  calculation shows that 
        $
        a < \frac 5{7\lambda} b
        $. 
        We now compute 
        \begin{align*}
            D &\geq \dist_X(y,p_1(0)) - \dist_X(x,p_1(0)) \\
            &\geq \frac 1 \lambda b - a - c \\
            &\geq \frac 2 {7\lambda }b - c \\
            &= \frac {2b}{7\lambda a + 7\lambda b}(a+b) - c \\
            &= \frac {1}{6\lambda}(a+b) - c \\
            &> D.
        \end{align*}
        This is clearly a contradiction, and so $\dist_X(x,y) \geq \frac 1 {8\lambda} \dist_Q(x,y) -c$.
        
        Finally, let us assume that $x \in p_1$ and $y \in \alpha$. Again, to ease calculation let us now parametrise $p_1$ and $\alpha$ so that $\alpha(0) = p_1(0)$. Write $x = p_1(a)$, $y = \alpha(b)$, and $D = \dist_X(x,y)$. We split into two cases. Suppose first that $a \geq \frac b{2\lambda}$. Then, since $p_1$ is a shortest path to $T_n$ (and thus $\alpha)$, we easily see that 
        $$
        D  \geq a = \frac{\lambda a}{3\lambda a} + \frac{2\lambda a}{3\lambda a} \geq \frac{1}{3\lambda}(a+b).
        $$
        Suppose now that $a < \frac{b}{2\lambda}$. More routine calculation shows that
        \begin{align*}
            D &\geq \dist_X(y,p_1(0)) - \dist_X(p_1(0),x) \\
            & \geq \frac1 \lambda b - c- a \\
            & \geq \frac 1 {3\lambda}(a+b) -c.
        \end{align*}
        We have considered all possible cases, whence the claim.
    \end{proof}

\end{document}
